\documentclass{jsarticle}
\usepackage{amsmath}
\usepackage{amssymb}
\begin{document}
\title{}
\author{}
\date{}
\maketitle

  オイラーの公式 アーベル多様体 母関数 極限値 組み合わせ多様体 
  座標変換 場の変換則 素粒子方程式 重力・反重力方程式
  エントロピー不変量 Higgs場 時間 質量 相対性原理 スピンネットワーク
  特異点理論 素数分布論 Hilbert空間 対称性 正規部分郡
  Thurston多様体 閉３次元空間 電弱相互理論 反重力の双対性
  電強相互理論 共時性 不確定性原理 確率論 宇宙と異次元空間
  一次独立 波動・粒子の相補性 ラプラス方程式 大域的微分方程式
  Zeta関数 単体量 微分幾何構造 量子群 カタストロフィー理論
  形態形成場 フロベニウスの定理 
  ハイゼンベルク方程式 大域的微分方程式の真逆の方程式
  宇宙の逆さまは原子 重力波 単体量における大きさの量 
  宇宙は空間と時間が無限から有限へと生成されたときから始まる
  遷移元素 ランタノイド アクチノイド
  原子 グループ毎に分けられる 
  性質 要素 が素数の束の集合になっている 
  時間の一方向性 
  質量 体積
  boson fermison

  素数の束がゼータ関数の素数分布論を形成していて、DATABASEとして
  Hilbert空間にスピンネットワークが貼りめぐらされている。このHilbert空間
  がThurston多様体の微分幾何構造を形成していて、正規部分群を対称性
  をパリティの破れより、非対称性を帯びて、ラプラス方程式がこの非対称性
  を大域的微分方程式により形づくられて、重力波が単体量の大きさでゼータ関数
  を表している。ハイゼンベルク方程式から原子を求められて、その値が微細構造定数
  から宇宙の逆さまになっている。宇宙は空間と時間が無限から有限へと生成された
  ときから始まっている。原子は光量子仮説で遷移元素からランタノイドとアクチノイド
  へとゼータ関数からグループ毎に性質と要素の素数の束の集合論になっていて、
  ５種類の力から、電磁力と弱い力と重力が合わさって電弱相互理論を統合されていて、
  反重力を入れると電磁力と強い力と合わすと電強相互理論にも統合されている。
  この電弱相互理論は時間の一方向性を性質として持っていて、反重力の双対性
  から電強相互理論をも時間の真逆の一方向性を帯びていて、ベクトルをこれらの
  力から量子力学を持ち出すと共時性から時間が止まっていることがこれらの理論
  から証明できる。
  Higgs場から光速度不変の法則を密度エネルギーからコントロールできていて、
  質量と時間が同じ地平線を世界線から同一視できて、トポロジーからこれらが
  同値ともわかる。この同値からスピンネットワークが体積の幾何学的変化
  から時間で各次元を行き来しているのを質量を持ち出すと、重力場が時空
  の湾曲度から質量が分かるが、この質量は時間と同型となり、スピンネットワーク
  は微分幾何の加速度の組み合わせ多様体から質量を時間と場の変換則を使い
  スピンネットワークが質量と体積の密度変化の光速度不変の法則の書き換えで
  スピノール場の式に行き着く。アーベル多様体はゼータ関数と合わさると
  カテゴリー理論へと行き着く。ゼータ関数は素数の束の各組み合わせ多様体から
  宇宙から原子までのレベルでアーベル多様体を積空間から接続すると
  いろんな物理学の基底数を導ける事ができる。
  
 
  
  
\end{document}
